\documentclass[11pt,a4paper]{article}
% Shared preamble for papers in this repository.
% Papers pick this up via TEXINPUTS (set by tools/paper/build.py), so a paper needs only
%
%   \documentclass[11pt,a4paper]{article}
%   % Shared preamble for papers in this repository.
% Papers pick this up via TEXINPUTS (set by tools/paper/build.py), so a paper needs only
%
%   \documentclass[11pt,a4paper]{article}
%   % Shared preamble for papers in this repository.
% Papers pick this up via TEXINPUTS (set by tools/paper/build.py), so a paper needs only
%
%   \documentclass[11pt,a4paper]{article}
%   \input{preamble}
%
% Keep this conservative: only packages present in a default MiKTeX/TeX Live install.

\usepackage[T1]{fontenc}
\usepackage[utf8]{inputenc}
\usepackage{lmodern}
\usepackage[margin=1in]{geometry}
\usepackage{amsmath, amssymb, amsthm}
\usepackage{mathtools}
\usepackage[hidelinks]{hyperref}
\usepackage{url}
\usepackage{booktabs}
\usepackage{enumitem}
\usepackage{microtype}

% Theorem environments, numbered within sections.
\theoremstyle{plain}
\newtheorem{theorem}{Theorem}[section]
\newtheorem{proposition}[theorem]{Proposition}
\newtheorem{lemma}[theorem]{Lemma}
\newtheorem{corollary}[theorem]{Corollary}
\newtheorem{conjecture}[theorem]{Conjecture}
\theoremstyle{definition}
\newtheorem{definition}[theorem]{Definition}
\newtheorem{remark}[theorem]{Remark}
\newtheorem{observation}[theorem]{Observation}

% Graph-theory shorthands used across these papers.
% NB: \deg is already a LaTeX primitive, so define a distinct name rather than
% redefining it -- \newcommand on an existing command is a fatal error.
\newcommand{\olddeg}{\operatorname{deg}}
\newcommand{\Vthree}{V_{3}}
\newcommand{\Vbig}{V_{\ge 4}}
\newcommand{\cardV}[1]{\lvert #1 \rvert}
\newcommand{\Ccyc}[1]{C_{#1}}

% A visible marker for statements that are machine-checked, with the Lean name.
\newcommand{\leanname}[1]{\textup{\texttt{#1}}}
\newcommand{\checkedin}[1]{\hfill\textup{[Lean: \leanname{#1}]}}

% Let bibliographies cite a DOI without pulling in another package.
\providecommand{\doi}[1]{\textsc{doi:}\,\url{#1}}

%
% Keep this conservative: only packages present in a default MiKTeX/TeX Live install.

\usepackage[T1]{fontenc}
\usepackage[utf8]{inputenc}
\usepackage{lmodern}
\usepackage[margin=1in]{geometry}
\usepackage{amsmath, amssymb, amsthm}
\usepackage{mathtools}
\usepackage[hidelinks]{hyperref}
\usepackage{url}
\usepackage{booktabs}
\usepackage{enumitem}
\usepackage{microtype}

% Theorem environments, numbered within sections.
\theoremstyle{plain}
\newtheorem{theorem}{Theorem}[section]
\newtheorem{proposition}[theorem]{Proposition}
\newtheorem{lemma}[theorem]{Lemma}
\newtheorem{corollary}[theorem]{Corollary}
\newtheorem{conjecture}[theorem]{Conjecture}
\theoremstyle{definition}
\newtheorem{definition}[theorem]{Definition}
\newtheorem{remark}[theorem]{Remark}
\newtheorem{observation}[theorem]{Observation}

% Graph-theory shorthands used across these papers.
% NB: \deg is already a LaTeX primitive, so define a distinct name rather than
% redefining it -- \newcommand on an existing command is a fatal error.
\newcommand{\olddeg}{\operatorname{deg}}
\newcommand{\Vthree}{V_{3}}
\newcommand{\Vbig}{V_{\ge 4}}
\newcommand{\cardV}[1]{\lvert #1 \rvert}
\newcommand{\Ccyc}[1]{C_{#1}}

% A visible marker for statements that are machine-checked, with the Lean name.
\newcommand{\leanname}[1]{\textup{\texttt{#1}}}
\newcommand{\checkedin}[1]{\hfill\textup{[Lean: \leanname{#1}]}}

% Let bibliographies cite a DOI without pulling in another package.
\providecommand{\doi}[1]{\textsc{doi:}\,\url{#1}}

%
% Keep this conservative: only packages present in a default MiKTeX/TeX Live install.

\usepackage[T1]{fontenc}
\usepackage[utf8]{inputenc}
\usepackage{lmodern}
\usepackage[margin=1in]{geometry}
\usepackage{amsmath, amssymb, amsthm}
\usepackage{mathtools}
\usepackage[hidelinks]{hyperref}
\usepackage{url}
\usepackage{booktabs}
\usepackage{enumitem}
\usepackage{microtype}

% Theorem environments, numbered within sections.
\theoremstyle{plain}
\newtheorem{theorem}{Theorem}[section]
\newtheorem{proposition}[theorem]{Proposition}
\newtheorem{lemma}[theorem]{Lemma}
\newtheorem{corollary}[theorem]{Corollary}
\newtheorem{conjecture}[theorem]{Conjecture}
\theoremstyle{definition}
\newtheorem{definition}[theorem]{Definition}
\newtheorem{remark}[theorem]{Remark}
\newtheorem{observation}[theorem]{Observation}

% Graph-theory shorthands used across these papers.
% NB: \deg is already a LaTeX primitive, so define a distinct name rather than
% redefining it -- \newcommand on an existing command is a fatal error.
\newcommand{\olddeg}{\operatorname{deg}}
\newcommand{\Vthree}{V_{3}}
\newcommand{\Vbig}{V_{\ge 4}}
\newcommand{\cardV}[1]{\lvert #1 \rvert}
\newcommand{\Ccyc}[1]{C_{#1}}

% A visible marker for statements that are machine-checked, with the Lean name.
\newcommand{\leanname}[1]{\textup{\texttt{#1}}}
\newcommand{\checkedin}[1]{\hfill\textup{[Lean: \leanname{#1}]}}

% Let bibliographies cite a DOI without pulling in another package.
\providecommand{\doi}[1]{\textsc{doi:}\,\url{#1}}


\title{The cubic density of a minimal counterexample to the\\
Erd\H{o}s--Gy\'arf\'as conjecture:\\
a reformulation, a machine-checked strict bound,\\
and an obstruction}

\author{Yuri Olive\thanks{Formalization and write-up. The strictness argument of
Section~\ref{sec:strict} is due to the Erd\H{o}s Problems forum user \texttt{jul059};
see Section~\ref{sec:attribution}.}}

\date{\today}

\begin{document}
\maketitle

\begin{abstract}
The Erd\H{o}s--Gy\'arf\'as conjecture asserts that every finite graph of minimum degree at
least $3$ contains a cycle whose length is a power of two. Carr proved that in a minimal
counterexample at least $4/7$ of the vertices are cubic; Bisch improved this to $2/3$. We
record three things. First, a reformulation: writing $S_3$ for the number of
cubic-to-cubic incidences, every minimal counterexample satisfies
$4\cardV{\Vbig} + S_3 \le 3\cardV{\Vthree}$, so the constant in any such density bound is
decided entirely by how many edges the cubic set spans among itself; Bisch's $2/3$ is
exactly the case in which $G[\Vthree]$ is a perfect matching. Second, a machine-checked
proof, in Lean~4 against \texttt{Mathlib}, of the strict bound
$\cardV{\Vthree} > \tfrac{2}{3}\cardV{V}$, an argument due to the forum user
\texttt{jul059} and previously unverified. Third, a conditional improvement: if
$G[\Vthree]$ has no $K_2$ component then $\cardV{\Vthree} \ge \tfrac{12}{17}\cardV{V}$,
together with an explanation of why removing that hypothesis resists the standard
minimality arguments, and an observation that no counting argument built on the known
structural lemmas can beat $2/3$. All statements below are formalized; nothing here is
asserted that the Lean development does not prove.
\end{abstract}

\section{Introduction}

\begin{conjecture}[Erd\H{o}s--Gy\'arf\'as]
Every finite graph with minimum degree at least $3$ contains a simple cycle whose length is
a power of two.
\end{conjecture}

The conjecture is open. Liu and Montgomery~\cite{LiuMontgomery} proved it for graphs whose
minimum degree exceeds an absolute constant, which also refuted the stronger conjecture of
Erd\H{o}s and Gy\'arf\'as; consequently a counterexample can only have very small minimum
degree. Markstr\"om~\cite{Markstrom} showed by computer search that a cubic counterexample
needs at least $30$ vertices, and Garcia~\cite{Garcia} recently established, by a
SAT search certified with DRAT proofs, that any counterexample has at least $24$ vertices.

Throughout, $G$ is a \emph{minimal counterexample}: a counterexample of minimum order and,
subject to that, of minimum size. Write
\[
  \Vthree = \{v : \olddeg v = 3\},
  \qquad
  \Vbig = \{v : \olddeg v \ge 4\}.
\]
Carr~\cite{Carr} proved that $\Vbig$ is an independent set, that every vertex has a
neighbour of degree exactly $3$, and that $\cardV{\Vthree} \ge \tfrac{4}{7}\cardV{V}$.
Bisch~\cite{Bisch} formalized those structural facts in Lean~4 and improved the density to
$\cardV{\Vthree} \ge \tfrac{2}{3}\cardV{V}$.

Our contribution is organised as follows. Section~\ref{sec:reform} gives the counting
reformulation, which makes the source of the constant explicit. Section~\ref{sec:strict}
reports the machine-checked strict bound. Section~\ref{sec:cond} proves the conditional
$12/17$ bound, and Section~\ref{sec:obstruction} explains why we could not remove its
hypothesis --- including the observation that counting alone cannot.

\paragraph{Standing hypotheses.}
All results are stated for a graph satisfying the properties of a minimal counterexample
that we actually use: minimum degree at least $3$; $\Vbig$ independent; every vertex has a
cubic neighbour; no $4$-cycle; minimality in the order; and the absence of a power-of-two
cycle. In the formalization these are bundled as \leanname{MinCexHyps}, and
\leanname{EGCBridge} derives that bundle from Bisch's \leanname{IsMinCex}, so every
statement below also holds verbatim for a minimal counterexample in his sense.

\section{The reformulation}\label{sec:reform}

\begin{definition}
For a vertex $v$ let $\operatorname{cd}(v)$ denote its number of cubic neighbours, and put
\[
  S_3 \;=\; \sum_{v \in \Vthree} \operatorname{cd}(v),
\]
the number of ordered cubic-to-cubic incidences; equivalently $S_3 = 2\,e(G[\Vthree])$.
\end{definition}

\begin{theorem}\label{thm:reform}
$4\cardV{\Vbig} + S_3 \le 3\cardV{\Vthree}$. \checkedin{four\_card\_big\_add\_sum\_cubicDeg\_le}
\end{theorem}

\begin{proof}
Since $\Vbig$ is independent, every neighbour of a vertex of $\Vbig$ is cubic, so
$\olddeg u$ counts cubic neighbours for each $u \in \Vbig$ and
$\sum_{u \in \Vbig} \olddeg u = e(\Vbig, \Vthree)$. Counting that quantity from the other
side and using that a cubic vertex has exactly $3 - \operatorname{cd}(v)$ neighbours in
$\Vbig$,
\[
  4\cardV{\Vbig}
  \;\le\; \sum_{u \in \Vbig} \olddeg u
  \;=\; e(\Vbig, \Vthree)
  \;=\; \sum_{v \in \Vthree}\bigl(3 - \operatorname{cd}(v)\bigr)
  \;=\; 3\cardV{\Vthree} - S_3. \qedhere
\]
\end{proof}

Theorem~\ref{thm:reform} isolates what controls the constant: a density bound follows from
any lower bound on $S_3$, and from nothing else.

\begin{corollary}[Bisch's bound]\label{cor:bisch}
$4\cardV{\Vbig} \le 2\cardV{\Vthree}$, hence
$\cardV{\Vthree} \ge \tfrac{2}{3}\cardV{V}$. \checkedin{card\_big\_le\_of\_domination}
\end{corollary}

\begin{proof}
Every cubic vertex has a cubic neighbour, so $\operatorname{cd}(v) \ge 1$ on $\Vthree$ and
$S_3 \ge \cardV{\Vthree}$. Substitute into Theorem~\ref{thm:reform}.
\end{proof}

\begin{observation}\label{obs:matching}
Equality $S_3 = \cardV{\Vthree}$ holds precisely when $G[\Vthree]$ is a perfect matching.
So $2/3$ is exactly the bound obtained in that case, and a $K_2$ component of
$G[\Vthree]$ --- a pair of adjacent cubic vertices whose four remaining edges all leave to
$\Vbig$ --- is the unique local configuration attaining the extremal ratio.
\end{observation}

\section{The strict bound, verified}\label{sec:strict}

\begin{theorem}\label{thm:strict}
$\cardV{\Vthree} \ge 2\cardV{\Vbig} + 1$, hence
$3\cardV{\Vthree} > 2\cardV{V}$. \checkedin{strict\_two\_thirds}
\end{theorem}

The argument is not ours; it is due to \texttt{jul059} \cite{jul059}, who posted it on the
Erd\H{o}s Problems forum marked as unverified. We state it for completeness and because the
formalization is what we contribute.

Suppose $\cardV{\Vthree} = 2\cardV{\Vbig}$. Then every inequality in the proof of
Theorem~\ref{thm:reform} is tight, which forces every vertex of $\Vbig$ to have degree
exactly $4$ and every cubic vertex to have exactly two neighbours in $\Vbig$
(\leanname{degree\_eq\_four\_of\_equality}, \leanname{card\_big\_nbrs\_eq\_two\_of\_equality}).
Contract each cubic vertex to an edge joining its two $\Vbig$ neighbours. The result $H$ is
a graph on $\Vbig$; it is simple because two cubic vertices with the same pair of $\Vbig$
neighbours would close a $4$-cycle, and it is $4$-regular because the four cubic
neighbours of a vertex $u \in \Vbig$ lead to four \emph{distinct} vertices --- a collision
would again close a $4$-cycle through $u$ (\leanname{contract\_degree\_ge}). Since
$\cardV{V(H)} = \cardV{\Vbig} < \cardV{V}$ and $H$ has minimum degree $4$, minimality
supplies a cycle in $H$ of length $2^k$. Re-inserting the cubic vertices turns it into a
cycle of length $2^{k+1}$ in $G$, contradicting the hypothesis on $G$.

\begin{remark}
The lifting step is where the formalization does work the informal argument leaves
implicit: one must check that the resulting closed walk is a \emph{cycle}. The $\Vbig$
vertices are distinct because the cycle in $H$ is one; the inserted cubic vertices are
distinct because, in the equality case, a cubic vertex has exactly two $\Vbig$ neighbours
and therefore \emph{determines} the contraction edge it came from
(\leanname{mid\_determines\_pair}), and a cycle traverses each edge once; and the two
families are disjoint because a cubic vertex has degree $3$ while a $\Vbig$ vertex does
not. See \leanname{exists\_pow2\_cycle\_of\_contract\_cycle}.
\end{remark}

\section{A conditional improvement}\label{sec:cond}

By Observation~\ref{obs:matching}, improving on $2/3$ means bounding $S_3$ away from
$\cardV{\Vthree}$, i.e.\ excluding $K_2$ components of $G[\Vthree]$.

\begin{definition}
$G$ has \emph{no $K_2$ component} if every cubic vertex with exactly one cubic neighbour
has that neighbour possess at least two. \leanname{noK2Component}
\end{definition}

\begin{lemma}\label{lem:fibre}
If $G$ has no $K_2$ component then $3 S_3 \ge 4\cardV{\Vthree}$.
\checkedin{three\_mul\_sum\_cubicDeg\_ge}
\end{lemma}

\begin{proof}
Split $\Vthree$ into $T = \{v : \operatorname{cd}(v) = 1\}$ and
$S = \{v : \operatorname{cd}(v) \ge 2\}$; domination excludes $\operatorname{cd}(v) = 0$,
so this is a partition and $S_3 = \cardV{T} + \sum_{s \in S}\operatorname{cd}(s)$. Write
$D = \sum_{s\in S}\operatorname{cd}(s)$, so $D \ge 2\cardV{S}$. The hypothesis sends each
$v \in T$ to a cubic neighbour lying in $S$; the fibre of $s \in S$ consists of cubic
neighbours of $s$, so it has at most $\operatorname{cd}(s)$ elements and hence
$\cardV{T} \le D$. Now
\[
  S_3 = \cardV{T} + D \ge \cardV{T} + \max\bigl(2\cardV{S},\, \cardV{T}\bigr),
\]
and comparing the two cases $\cardV{T} \ge 2\cardV{S}$ and $\cardV{T} < 2\cardV{S}$ gives
$S_3 \ge \tfrac{4}{3}\bigl(\cardV{T} + \cardV{S}\bigr) = \tfrac{4}{3}\cardV{\Vthree}$
in both.
\end{proof}

Note that the proof is a fibre count; no connectivity or spanning-tree argument is needed.

\begin{theorem}\label{thm:1217}
If $G$ has no $K_2$ component then $12\cardV{V} \le 17\cardV{\Vthree}$, i.e.
\[
  \cardV{\Vthree} \;\ge\; \tfrac{12}{17}\cardV{V} \;\approx\; 0.7059\,\cardV{V}.
\]
\checkedin{twelve\_seventeenths\_of\_noK2}
\end{theorem}

\begin{proof}
Combine Theorem~\ref{thm:reform} with Lemma~\ref{lem:fibre}:
$4\cardV{\Vbig} \le 3\cardV{\Vthree} - \tfrac{4}{3}\cardV{\Vthree} = \tfrac{5}{3}\cardV{\Vthree}$,
so $12\cardV{\Vbig} \le 5\cardV{\Vthree}$. Adding $12\cardV{\Vthree}$ to both sides and
using $\cardV{V} = \cardV{\Vthree} + \cardV{\Vbig}$ gives the claim.
\end{proof}

\section{Why the hypothesis resists removal}\label{sec:obstruction}

Removing the hypothesis of Theorem~\ref{thm:1217} would give $12/17$ outright. We did not
manage it, and we record two obstacles, both of which we checked, so that the route is not
re-attempted blindly.

\subsection*{Local replacements shift cycle lengths by one}

Bisch~\cite{Bisch} constrains the offending configuration tightly: if $v, v'$ are adjacent
and every other neighbour of either has degree at least $4$, then $v$ and $v'$ have a
\emph{unique} common neighbour $u$, and $\olddeg u = 4$. So each $K_2$ component comes
with a triangle $v v' u$ whose apex has degree exactly $4$, its two remaining neighbours
being cubic.

Every way we found of destroying such a component changes cycle lengths additively:

\begin{itemize}[nosep]
  \item delete $v, v'$ and join $u$ to the two outer $\Vbig$ neighbours: a $2^k$ cycle
        upstairs lifts to $2^k + 1$ or $2^k + 2$;
  \item delete $v, v', u$ and join the apex's other two neighbours: lifts to $2^k+1$;
  \item contract the triangle to a single vertex: shortens cycles by $0$ or $1$ depending
        on the entry and exit edges, hence non-uniformly;
  \item delete the apex and rewire its neighbours in pairs: lifts to $2^k+1$.
\end{itemize}

Since $2^k + 1$ and $2^k+2$ are not powers of two, minimality yields no contradiction in
any of these cases. The obstruction is structural rather than a failure of ingenuity:
replacing a $2$-path by an edge shifts every cycle through it by exactly one, and the
powers of two are too sparse for an additive shift to land on another one. The argument of
Section~\ref{sec:strict} escapes this only because in the perfect-matching case the
operation is \emph{uniformly multiplicative} --- $G$ is a subdivision of the contraction,
every cycle doubles, and $2^k \mapsto 2^{k+1}$. A single edge inside $\Vthree$ destroys
that uniformity.

\subsection*{Counting alone cannot beat two thirds}

More is true: the extremal configuration appears to be consistent with all of the
structural facts currently available. Take $G[\Vthree]$ to be a perfect matching in which
every degree-$4$ apex serves two $K_2$ components --- so $u \sim v, v', p, p'$ with
$v \sim v'$ and $p \sim p'$ --- and let the remaining $\Vbig$ vertices have degree $4$,
each absorbing one edge from four distinct components. Then $\Vbig$ is independent, every
vertex has a cubic neighbour, the uniqueness in Bisch's proposition is respected, no
$4$-cycle is forced, and the count closes at exactly $\cardV{\Vbig} = \cardV{\Vthree}/2$.

Consequently no argument that uses only Theorem~\ref{thm:reform} together with the known
structural lemmas can improve on $2/3$: any improvement must invoke the cycle condition,
which is precisely what the additive obstruction above blocks. Breaking the deadlock needs
either an operation on the $K_2$ configuration that scales all cycle lengths
multiplicatively, or a new structural lemma about the apexes.

\section{Formalization}\label{sec:formal}

Every numbered statement above is proved in Lean~4 against \texttt{Mathlib}, with no use of
\texttt{sorry}. The development is organised as

\begin{center}
\begin{tabular}{ll}
\toprule
File & Contents \\
\midrule
\leanname{EGCStrict.lean}  & hypothesis bundle, counting, the equality analysis, the contraction \\
\leanname{EGCLift.lean}    & the cycle lifting; Theorem~\ref{thm:strict} \\
\leanname{EGCDensity.lean} & Theorem~\ref{thm:reform}, Corollary~\ref{cor:bisch}, Theorem~\ref{thm:1217} \\
\leanname{EGCBridge.lean}  & derives the hypothesis bundle from Bisch's \leanname{IsMinCex} \\
\bottomrule
\end{tabular}
\end{center}

Each theorem was checked with \verb|#print axioms| to depend only on \texttt{propext},
\texttt{Classical.choice} and \texttt{Quot.sound}; the audit files are
\leanname{AxiomCheck.lean}, \leanname{BridgeCheck.lean} and \leanname{DensityCheck.lean}.
Bisch's file is not redistributed --- its repository carries no licence --- so
\leanname{EGCBridge} is not a default build target and expects the file to be fetched
locally; the other three modules build and prove their statements without it, conditional
on the hypothesis bundle.

\section{Attribution}\label{sec:attribution}

The mathematics of Section~\ref{sec:strict} is due to the forum user \texttt{jul059}
\cite{jul059}. Their comment supplies the equality analysis, the contraction, the reason it
is simple, and the doubling of cycle lengths; it was posted explicitly as unverified. Our
contribution there is the formalization, including the distinctness bookkeeping. The
structural lemmas are Carr's~\cite{Carr}, formalized by Bisch~\cite{Bisch}, who also
proved the non-strict $2/3$ bound. Sections~\ref{sec:reform}, \ref{sec:cond}
and~\ref{sec:obstruction} are ours.

We would welcome direction from \texttt{jul059} on how they wish to be credited; we have
used the handle because no name is attached to the posting.

\begin{thebibliography}{9}

\bibitem{Bisch}
A.~Bisch.
\newblock \emph{On the density of cubic vertices in a minimal counterexample to the
Erd\H{o}s--Gy\'arf\'as conjecture}.
\newblock Zenodo, \doi{10.5281/zenodo.21574476}. Lean formalization:
\url{https://github.com/AJBisch/AJBisch.github.io/blob/main/EGC.lean}.

\bibitem{Carr}
A.~Carr.
\newblock \emph{Every minimal counterexample to the Erd\H{o}s--Gy\'arf\'as conjecture is
predominantly cubic}.
\newblock arXiv:2605.22844.

\bibitem{Garcia}
D.~Garcia.
\newblock \emph{Small graphs without power-of-two cycles: a lower bound of 24, a correction
to a construction of Exoo, and explicit bounds for $f(k)$}.
\newblock arXiv:2609.04686.

\bibitem{jul059}
\texttt{jul059}.
\newblock Comment on Erd\H{o}s Problem \#64, 26 July 2026.
\newblock \url{https://www.erdosproblems.com/forum/thread/64\#post-8130}.

\bibitem{LiuMontgomery}
H.~Liu and R.~Montgomery.
\newblock \emph{A solution to Erd\H{o}s and Hajnal's odd cycle problem}, and related work on
cycle lengths in graphs of large average degree.
\newblock See the discussion at \url{https://www.erdosproblems.com/64}.

\bibitem{Markstrom}
K.~Markstr\"om.
\newblock \emph{Extremal graphs for some problems on cycles in graphs}.
\newblock Congressus Numerantium 171 (2004).

\end{thebibliography}

\end{document}
